\begin{solapartat}
El camp elèctric és una magnitud vectorial per tant:
\[ \vec{E}_T = \vec{E}_1 + \vec{E}_2 = K\,\frac{q_1\,\vec{\mu}_1}{r_1^2} + K\,\frac{q_2\,\vec{\mu}_2}{r_2^2} \quad \text{\punts{0,2}} \]
On:
\[ \vec{\mu}_1 = \left(\frac{\vec{\imath}}{\sqrt{2}} - \frac{\vec{\jmath}}{\sqrt{2}}\right);\quad \vec{\mu}_2 = \left(\frac{\vec{\imath}}{\sqrt{2}} + \frac{\vec{\jmath}}{\sqrt{2}}\right);\quad r_1 = r_2 = \sqrt{2} \quad \text{\punts{0,3}} \]
Per tant:
\begin{align*}
&\vec{E}_T = \frac{9\cdot 10^{9}}{2\,\sqrt{2}}\left\{9\cdot 10^{-6}(\vec{\imath} - \vec{\jmath}) - 9\cdot 10^{-6}(\vec{\imath} + \vec{\jmath})\right\}\\
&\Rightarrow \vec{E}_T = (0\,\vec{\imath} - 5,73\cdot 10^{4}\,\vec{\jmath})\ \mathrm{N/C} \quad \text{\punts{0,5}}
\end{align*}
Es considerarà la resposta correcte si raonen que per raons de simetria el camp elèctric ha de tenir només component vertical, de signe negatiu i realitzen el càlcul correctament.
\end{solapartat}

\begin{solapartat}
Al tractar-se de un camp conservatiu podem podem trobar el treball fet per la força elèctrica a partir del potencial elèctric.
\[ V_i = K\left\{\frac{q_1}{r_1} + \frac{q_2}{r_2}\right\} = \frac{9\cdot 10^{9}\cdot 10^{-6}}{\sqrt{2}}\,(9 - 9) = 0\ \mathrm{V} \quad \text{\punts{0,3}} \]
\[ V_f = K\left\{\frac{q_1}{r'_1} + \frac{q_2}{r'_2}\right\} = 9\cdot 10^{9}\cdot 10^{-6}\left(\frac{9}{2\sqrt{2}} - \frac{9}{2}\right) = -1,19\cdot 10^{4}\ \mathrm{V} \quad \text{\punts{0,3}} \]
Per tant el treball fet per la força elèctrica és:
\[ W_E = -\Delta V\,q = (V_i - V_f)\,q = (0 + 1,19\cdot 10^{4})\,(7\cdot 10^{-6}) = 8,33\cdot 10^{-2}\ \mathrm{J} \quad \text{\punts{0,4}} \]
\end{solapartat}
